% TODO make hw stylesheet
\documentclass[10pt]{article}
\usepackage[letterpaper,top=1in,left=1in,right=1in]{geometry}
\usepackage[dvipsnames,svgnames]{xcolor}
\usepackage[shortlabels]{enumitem}
\usepackage[framemethod=TikZ]{mdframed}
\usepackage{amsmath,amssymb,amsthm}
\usepackage[colorlinks]{hyperref}
\usepackage{multicol}
\usepackage{tabularx}
\usepackage{bbm}

\hypersetup{
    colorlinks=true,
    linkcolor=blue,
    filecolor=magenta,
    urlcolor=magenta,
    pdftitle={MAT 344 HW}
}

\usepackage{mathtools}
\usepackage{thmtools}
\usepackage[textsize=footnotesize]{todonotes}
\usepackage{titling}
\usepackage{cancel}

\reversemarginpar
\mdfdefinestyle{mdbluebox}{roundcorner=10pt,innerbottommargin=9pt,
linecolor=blue,backgroundcolor=TealBlue!5,}
\declaretheoremstyle[headfont=\sffamily\bfseries\color{MidnightBlue},
mdframed={style=mdbluebox},]{thmbluebox}

\mdfdefinestyle{mdbloobox}{frametitlefont=\bfseries,innerbottommargin=8pt,
nobreak=true,backgroundcolor=TealBlue!5,linecolor=blue,}
\declaretheoremstyle[headfont=\bfseries\color{MidnightBlue},
mdframed={style=mdbloobox},headpunct={\\[3pt]},postheadspace=0pt,]{thmbloobox}

\mdfdefinestyle{mdredbox}{frametitlefont=\bfseries,innerbottommargin=8pt,
nobreak=true,backgroundcolor=Salmon!5,linecolor=RawSienna,}
\declaretheoremstyle[headfont=\bfseries\color{RawSienna},
mdframed={style=mdredbox},headpunct={\\[3pt]},postheadspace=0pt,]{thmredbox}

\mdfdefinestyle{mdgreenbox}{linecolor=ForestGreen,backgroundcolor=ForestGreen!5,
linewidth=2pt,rightline=false,leftline=true,topline=false,bottomline=false,}
\declaretheoremstyle[headfont=\bfseries\sffamily\color{ForestGreen!70!black},
mdframed={style=mdgreenbox},headpunct={ --- },]{thmgreenbox}

\mdfdefinestyle{mdorangebox}{linecolor=RedOrange,backgroundcolor=RedOrange!5,
linewidth=2pt,rightline=false,leftline=true,topline=false,bottomline=false,}
\declaretheoremstyle[headfont=\bfseries\sffamily\color{RedOrange!70!black},
mdframed={style=mdorangebox},headpunct={ --- },]{thmorangebox}

\mdfdefinestyle{mdblackbox}{linecolor=black,backgroundcolor=RedViolet!5!gray!5,
linewidth=3pt,rightline=false,leftline=true,topline=false,bottomline=false,}
\declaretheoremstyle[mdframed={style=mdblackbox}]{thmblackbox}

\declaretheoremstyle[spaceabove=6pt,spacebelow=6pt]{boldhead}
\declaretheoremstyle[spaceabove=6pt,spacebelow=6pt,headfont=\normalfont\itshape]{italichead}

\declaretheorem[style=boldhead,name=Problem]{problem}
\declaretheorem[style=boldhead,name=Problem,numbered=no]{problem*}
\declaretheorem[style=boldhead,name=Setup]{setup} % for config geo
\declaretheorem[style=boldhead,name=Example,sibling=problem]{example}
\declaretheorem[style=boldhead,name=Theorem,sibling=problem]{theorem}
\declaretheorem[style=boldhead,name=Corollary,sibling=theorem]{corollary}
\declaretheorem[style=boldhead,name=Proposition,sibling=theorem]{prop}
\declaretheorem[style=boldhead,name=Lemma,numberwithin=problem]{lemma}
\declaretheorem[style=boldhead,name=Theorem,numbered=no]{theorem*}
\declaretheorem[style=boldhead,name=Lemma,numbered=no]{lemma*}
\declaretheorem[style=boldhead,name=Claim,numberwithin=problem]{claim}
\declaretheorem[style=boldhead,name=Claim,numbered=no]{claim*}
\declaretheorem[style=boldhead,name=Remark,numberwithin=problem]{remark}
\declaretheorem[style=boldhead,name=Remark,numbered=no]{remark*}
\declaretheorem[style=boldhead,name=Definition,sibling=theorem]{definition}
\declaretheorem[style=boldhead,name=Definition,numbered=no]{definition*}

\providecommand{\ol}{\overline}
\providecommand{\ceil}[1]{\left\lceil #1 \right\rceil}
\providecommand{\floor}[1]{\left\lfloor #1 \right\rfloor}
\newcommand{\dd}{\mathrm{d}}
\providecommand{\ul}{\underline}
\providecommand{\eps}{\varepsilon}
\providecommand{\half}{\frac{1}{2}}
\providecommand{\CC}{\mathbb C}
\providecommand{\FF}{\mathbb F}
\providecommand{\II}{\mathbb I}
\providecommand{\NN}{\mathbb N}
\providecommand{\QQ}{\mathbb Q}
\providecommand{\RR}{\mathbb R}
\providecommand{\ZZ}{\mathbb Z}
\providecommand{\ind}{\mathbbm 1}
\providecommand{\dg}{^\circ}
\providecommand{\ii}{\item}
\providecommand{\setdiff}{\smallsetminus}
\providecommand{\alert}[1]{{\sffamily\textbf{\textcolor{blue}{#1}}}}
\providecommand{\inv}{^{-1}}
\providecommand{\mb}[1]{\mathbf{#1}}
\newcommand{\what}{\widehat}

\newenvironment{soln}{\begin{proof}[Solution]}{\end{proof}}

\providecommand{\pow}{\operatorname{Pow}}
\providecommand{\vocab}[1]{{\textbf{\textcolor{ForestGreen}{#1}}}}
\providecommand{\lcm}{\operatorname{lcm}}
\providecommand{\abs}[1]{\left\lvert #1\right\rvert}
\providecommand{\smabs}[1]{\lvert #1\rvert}
\providecommand{\innerprod}[1]{\langle\!\langle #1\rangle\!\rangle}
\providecommand{\norm}[1]{\left\| #1\right\|}
\providecommand{\smnorm}[1]{\| #1\|}
\providecommand{\gr}{\operatorname{gr}}

\usepackage{mathrsfs}

% place title at top right
\setlength{\droptitle}{-8ex}
\pretitle{\begin{flushright}\Large\bfseries}
\posttitle{\par\end{flushright}}
\preauthor{\begin{flushright}}
\postauthor{\end{flushright}}
\predate{\begin{flushright}}
\postdate{\end{flushright}}

\title{MAT 344 HW02}
\author{\large Aditya Pahuja}
\date{\large Due 02/10}
\begin{document}
\maketitle

\begin{problem}
    Show for a complex valued function $f$ that
    \begin{equation*}
	\dd f = \frac{\partial f}{\partial z}\dd z
	+ \frac{\partial f}{\partial \ol z}\dd\ol z.
    \end{equation*}
\end{problem}

\begin{soln}
    Recalling that
    \begin{equation*}
	\frac{\partial f}{\partial z}
	= \half\left(\frac{\partial f}{\partial x}
	- i\frac{\partial f}{\partial y}\right),\qquad
	\frac{\partial f}{\partial z}
	= \half\left(\frac{\partial f}{\partial x}
	+ i\frac{\partial f}{\partial y}\right),
    \end{equation*}
    we directly compute
    \begin{align*}
        \frac{\partial f}{\partial z}\dd z
	+ \frac{\partial f}{\partial \ol z}\dd z
	&= \half\left(\frac{\partial f}{\partial x}
	- i\frac{\partial f}{\partial y}\right)\cdot
	(\dd x + i\dd y)
	+ \half\left(\frac{\partial f}{\partial x}
	+ i\frac{\partial f}{\partial y}\right)\cdot
	(\dd x - i\dd y)\\
	&= \frac{\partial f}{\partial f}\dd x
	+ \frac{\partial f}{\partial y}\dd y = \dd f.\qedhere
    \end{align*}
\end{soln}

\begin{problem}
    For which values of $z$ is
    \begin{equation*}
	\sum_{n=0}^\infty \left(\frac z{1+z}\right)^n
    \end{equation*}
    convergent?
\end{problem}

\begin{soln}
    A geometric series of complex numbers converges
    if and only if the common ratio of the sequence being summed
    has magnitude smaller than $1$, so we want to find
    the subset of $\CC$ containing all $z$ such that
    $\smabs{\frac{z}{1+z}} < 1$. Letting $z = a+bi$ for real $a,b$,
    this is equivalent to
    \begin{equation*}
	\sqrt{a^2+b^2} < \sqrt{(1+a)^2 + b^2}
	\iff a^2 < (1+a)^2 \iff -\half < a,
    \end{equation*}
    i.e. the series converges for \framebox{all $z$ such that
    $\operatorname{Re}(z) > -\half$}.
\end{soln}

\begin{problem}
    Using the skew-symmetric rules for wedge multiplication
    of real $1$-forms, show that
    \begin{equation*}
        \dd z\wedge\dd z = \dd\ol z\wedge\dd\ol z = 0
	\qquad\text{and}\qquad
	\dd z\wedge\dd\ol z = -\dd\ol z\wedge\dd z
	= -2i\dd x\wedge\dd y.
    \end{equation*}
    Use this to express the formula in Green's theorem
    for $\omega = f(z)\dd z$
    (where $f(x,y) = u(x,y) + iv(x,y)$) in terms of
    $\frac{\partial f}{\partial\ol z}\dd z\wedge\dd\ol z$.
\end{problem}

\begin{soln}
    Wedging a $1$-form with itself is always zero because
    swapping the arguments yields
    $\omega\wedge\omega = -\omega\wedge\omega$
    by anticommutativity.
    This implies that $2\omega\wedge\omega = \omega\wedge\omega = 0$,
    so in particular $\dd z\wedge\dd z = \dd\ol z\wedge\dd\ol z = 0$.
    The other two equations follow from a direct calculation:
    \begin{gather*}
	\dd z\wedge\dd\ol z
	= (\dd x + i\dd y)(\dd x - i\dd y)
	= \dd x\wedge\dd x + \dd y\wedge\dd y
	+ i(\dd y\wedge\dd x - \dd x\wedge\dd y)
	= -2i\dd x\wedge\dd y,
    \end{gather*}
    and this is equal to $-\dd\ol z\wedge\dd z$ by anticommutativity.
    Green's theorem says that, for a good domain $U$,
    \begin{equation*}
	\int_{\partial U}\omega = \int_U\dd\omega
	= \int_U \dd f\wedge\dd z
	= \int_U \left(\frac{\partial f}{\partial z}\dd z
	+ \frac{\partial f}{\partial\ol z}\dd\ol z\right)\wedge\dd z
	= \int_U \frac{\partial f}{\partial\ol z}\dd z\wedge\dd\ol z,
    \end{equation*}
    where the other summand in the integral cancels out
    because $\dd z\wedge\dd z = 0$.
\end{soln}

\pagebreak
\begin{problem}
    Let $\Delta = \{z\in\CC : \abs z\le 1\}$
    and choose $a\in\operatorname{Int}\Delta$.
    Let $\eps>0$ be small enough that
    $\Delta(\eps,a) \equiv \{z\in\CC : \abs{z-a} < \eps\}\subseteq
    \operatorname{Int}\Delta$. Let $f$ be a holomorphic function
    in a neighborhood of $\Delta$. Write out Green's formula for
    the holomorphic function
    \begin{equation*}
	F(z) = \frac {f(z)}{z-a}
    \end{equation*}
    on $\Omega \equiv\Delta - \operatorname{Int}\Delta(\eps,a)$.
    What happens as $\eps\to 0$?
\end{problem}

\begin{soln}
    Let $g(z) = \frac{f(z)}{z-a}$. Since $a\notin\Omega$,
    $\frac 1{z-a}$ is holomorphic on $\Omega$, hence so is $g$,
    being the product of two functions that are holomorphic on $\Omega$.
    In particular, $\frac{\partial g}{\partial\ol z} = 0$ on $\Omega$, so
    \begin{equation*}
	\int_{\partial\Omega}g(z)\dd z
	= \int_\Omega \frac{\partial g}{\partial\ol z}
	\dd z\wedge\dd\ol z = 0
    \end{equation*}
    by the result of Problem 3.
    If we orient the boundaries of $\Delta$ and $\Delta(\eps,a)$
    counterclockwise, the left-hand side equalling zero implies
    \begin{equation*}
	\int_{\partial\Delta} g(z)\dd z
	= \int_{\partial\Delta(\eps,a)}g(z)\dd z.\tag{$*$}
    \end{equation*}
    Let $\gamma_r = \partial\Delta(r,a)$,
    which we can parametrize as the curve
    $\gamma_r\colon[0,2\pi]\to\CC$ given by
    $\gamma_r(\theta) = a + re^{i\theta}$. Then,
    \begin{equation*}
	\int_{\gamma_r} \frac{f(z)}{z-a}\dd z
	= \int_0^{2\pi} \frac{f(a + r e^{i\theta})}
	{r e^{i\theta}}\cdot rie^{i\theta}\dd\theta
	= \int_0^{2\pi} if(a + re^{i\theta})\dd\theta.
    \end{equation*}
    Since $f$ is holomorphic at $a$, there exists a constant $C>0$
    such that, for each $r>0$, there is a $\delta>0$ satisfying
    $\abs{f(a+re^{i\theta}) - f(a)} < C\abs{re^{i\theta}} = Cr$
    (in particular, any $C > \abs{f'(a)}$ has this property).
    This implies that
    \begin{equation*}
	\abs{\int_0^{2\pi} if(a+re^{i\theta})\dd\theta
	- 2\pi if(a)}
	= \abs{\int_0^{2\pi} i(f(a + re^{i\theta}) - f(a))\dd\theta}
	< \int_0^{2\pi} Cr\dd\theta = 2\pi Cr,
    \end{equation*}
    so, in the limit,
    \begin{equation*}
	\lim_{r\to 0}\int_{\gamma_r} \frac{f(z)}{z-a}\dd z = f(a).
    \end{equation*}
    By $(*)$, the argument of the limit is constant as $r$ varies,
    so in fact this implies
    \begin{equation*}
	\int_{\partial\Delta}\frac{f(z)}{z-a}\dd z
	= \int_{\gamma_r}\frac{f(z)}{z-a}\dd z = f(a).\qedhere
    \end{equation*}
\end{soln}










\end{document}
